% Part: computability
% Chapter: recursive-functions
% Section: pr-functions

\documentclass[../../../include/open-logic-section]{subfiles}

\begin{document}

\olfileid{cmp}{rec}{prf}
\olsection{Fungsi Rekursif Primitif}

Ayo kita cathet maneh carane netepake fungsi anyar saka fungsi
kang wis ana nganggo rekursi primitif lan komposisi.

\begin{defn}
  \ollabel{defn:primitive-recursion} Upamana $f$ iku fungsi
  $k$ panggonan ($k\ge 1$) lan $g$ iku fungsi $(k+2)$ panggonan.
  Fungsi kang ditegesi kanthi \emph{rekursi primitif saka $f$ lan $g$}
  iku fungsi $(k+1)$ panggonan~$h$ kang ditemtokake dening
  persamaan iki:
  \begin{align*}
    h(x_0,\dots,x_{k-1},0) & =  f(x_0,\dots,x_{k-1}) \\
    h(x_0,\dots,x_{k-1},y+1) & = g(x_0,\dots,x_{k-1}, y, h(x_0,\dots,x_{k-1}, y))
  \end{align*}
\end{defn}

\begin{defn}
  \ollabel{defn:composition} Upamana $f$ iku fungsi $k$ panggonan,
  lan $g_0$, \dots, $g_{k-1}$ iku $k$ fungsi kang kabeh nduweni
  $n$ panggonan. Fungsi kang ditegesi kanthi \emph{komposisi saka $f$
    lan $g_0$, \dots,~$g_{k-1}$} iku fungsi $n$ panggonan~$h$
  kang ditemtokake dening
  \[
  h(x_0, \dots, x_{n-1}) =
  f(g_0(x_0, \dots, x_{n-1}), \dots, g_{k-1}(x_0, \dots, x_{n-1})).
  \]
\end{defn}

Saliyane $\Succ$ lan fungsi proyeksi
\[
\Proj{n}{i}(x_0,\dots,x_{n-1}) = x_i,
\]
kanggo saben wilangan asli $n$ lan $i < n$, kita uga bakal nglebokake
fungsi~$\Zero(x) = 0$ ing golongan fungsi rekursif primitif.

\begin{defn}
  Himpunan fungsi rekursif primitif iku himpunan fungsi
  saka $\Nat^n$ menyang $\Nat$, kang ditegesi kanthi induktif
  liwat syarat-syarat iki:
  \begin{enumerate}
  \item $\Zero$ iku rekursif primitif.
  \item $\Succ$ iku rekursif primitif.
  \item Saben fungsi proyeksi $\Proj{n}{i}$ iku rekursif primitif.
  \item Yen $f$ iku fungsi rekursif primitif $k$ panggonan lan $g_0$,
    \dots,~$g_{k-1}$ iku fungsi rekursif primitif $n$ panggonan,
    komposisi $f$ karo $g_0$, \dots,~$g_{k-1}$ uga rekursif
    primitif.
  \item Yen $f$ iku fungsi rekursif primitif $k$ panggonan lan $g$ iku
    fungsi rekursif primitif $k+2$ panggonan, fungsi kang ditegesi
    kanthi rekursi primitif saka $f$ lan $g$ uga rekursif
    primitif.
\end{enumerate}
\end{defn}

\begin{explain}
Kanthi luwih ringkes, himpunan fungsi rekursif primitif iku
himpunan paling cilik kang ngemot $\Zero$, $\Succ$, lan fungsi
proyeksi~$\Proj{n}{j}$, lan kang katutup tumrap komposisi lan
rekursi primitif.

Cara liyane kanggo nggambarake himpunan fungsi rekursif primitif
yaiku netepake miturut “tahap.” Upamana $S_0$ minangka himpunan
fungsi wiwitan: $\Zero$, $\Succ$, lan fungsi proyeksi. Iki
fungsi rekursif primitif tahap~$0$. Sawise tahap $S_i$ ditemtokake,
upamana $S_{i+1}$ minangka himpunan kabeh fungsi kang dipikolehi
kanthi ngetrapake siji langkah komposisi utawa rekursi primitif
marang fungsi-fungsi kang wis ana ing~$S_i$. Banjur
\[
S = \bigcup_{i \in \Nat} S_i
\]
iku himpunan kabeh fungsi rekursif primitif.
\end{explain}

Ayo dibuktekake yen $\Add$ iku fungsi rekursif primitif.

\begin{prop}
  Fungsi panjumlahan $\Add(x,y) = x+y$ iku rekursif primitif.
\end{prop}

\begin{proof}
Kita wis nduweni definisi rekursif primitif kanggo $\Add$ nganggo
rong fungsi $f$ lan~$g$ kang cocog karo pola ing
\olref{defn:primitive-recursion}:
\begin{align*}
  \Add(x_0, 0) & = f(x_0) = x_0\\
  \Add(x_0, y+1) & = g(x_0, y, \Add(x_0, y)) =
  \Succ(\Add(x_0, y))
\end{align*}
Dadi $\Add$ iku rekursif primitif yen $f$ lan $g$ uga
rekursif primitif. $f(x_0) = x_0 = \Proj{1}{0}(x_0)$, lan fungsi
proyeksi kalebu rekursif primitif, mula $f$ iku rekursif primitif.
Fungsi $g$ iku fungsi telung panggonan $g(x_0, y, z)$ kang
ditegesi kanthi
\[
g(x_0, y, z) = \Succ(z).
\]
Iki durung mbuktekake yen $g$ iku rekursif primitif, amarga $g$
lan $\Succ$ dudu fungsi kang persis padha: $\Succ$ iku fungsi
siji panggonan, dene $g$ kudu telung panggonan. Nanging kita
bisa netepake $g$ kanthi cara kang cocog karo aturan komposisi,
yaiku
\[
g(x_0, y, z) = \Succ(\Proj{3}{2}(x_0, y, z))
\]
Amarga $\Succ$ lan $\Proj{3}{2}$ kalebu fungsi rekursif
primitif, $g$ uga mangkono, awit bisa ditegesi kanthi komposisi
saka fungsi rekursif primitif.
\end{proof}

\begin{prop}
  \ollabel{prop:mult-pr}
  Fungsi perkalian $\Mult(x,y) = x \cdot y$ iku rekursif primitif.
\end{prop}

\begin{proof}
  Latihan.
\end{proof}

\begin{prob}
  Buktikna \olref[cmp][rec][prf]{prop:mult-pr} kanthi nuduhake yen
  definisi rekursif primitif kanggo $\Mult$ bisa diwujudake miturut
  pola kang diwajibake dening \olref[cmp][rec][prf]{defn:primitive-recursion}
  lan nuduhake yen fungsi $f$ lan~$g$ kang cocog iku rekursif
  primitif.
\end{prob}

\begin{ex}
Mangkene tuladha pisanan kita kanggo definisi rekursif primitif:
\begin{align*}
h(0) & =  1 \\
h(y+1) & =  2 \cdot h(y).
\end{align*}
Fungsi iki ora bisa mlebu pola kang diwajibake dening
\olref{defn:primitive-recursion}, amarga $k=0$. Definisi iki uga
nganggo konstanta $1$ lan $2$. Kanggo ngatasi masalah kapisan,
ayo kita tambahake argumen semu lan netepake fungsi~$h'$:
\begin{align*}
h'(x_0, 0) & =  f(x_0) = 1 \\
h'(x_0, y+1) & =  g(x_0, y, h'(x_0, y)) = 2 \cdot h'(x_0, y).
\end{align*}
Fungsi $f(x_0) = 1$ bisa ditegesi saka $\Succ$ lan $\Zero$
kanthi komposisi: $f(x_0) = \Succ(\Zero(x_0))$. Fungsi $g$ bisa
ditegesi kanthi komposisi saka $g'(z) = 2 \cdot z$ lan fungsi
proyeksi:
\begin{align*}
g(x_0, y, z) & = g'(\Proj{3}{2}(x_0, y, z))
\intertext{lan $g'$ banjur bisa ditegesi kanthi komposisi kaya mangkene}
g'(z) & = \Mult(g''(z), \Proj{1}{0}(z))
\intertext{lan}
g''(z) & = \Succ(f(z)),
\end{align*}
ing ngendi $f$ kaya ing ndhuwur: $f(z) = \Succ(\Zero(z))$.
Saiki, sawisé nduweni~$h'$, kita bisa nganggo komposisi maneh
kanggo netepake $h(y) =
h'(\Proj{1}{0}(y),\Proj{1}{0}(y))$. Iki nuduhake yen $h$ bisa
ditegesi saka fungsi dhasar liwat rerangken komposisi lan
rekursi primitif, mula $h$~iku rekursif primitif.
\end{ex}

\end{document}
